\section{Residue-relative weighted ideal counts}\label{sec:analytic}

Write
\[
 a_K(m)=\#\{\mathfrak a\subseteq R:\N\mathfrak a=m\},
 \qquad \kappaK=\operatorname*{Res}_{s=1}\zeta_K(s)>0.
\]
We first retain the residue in the analytic estimates. All constants
subscripted by $n$ are uniform over the fields under consideration.

\begin{lemma}\label{lem:residue}
For $D$ sufficiently large in terms of $n$,
\begin{equation}\label{eq:residue}
 h\zeta_K(1+h)\asymp_n\kappaK
 \quad\left(0<h\le\frac4{\log D}\right),
 \qquad \kappaK^{-1}\ll_n\log D.
\end{equation}
For the remaining fields, the first comparison holds uniformly for
sufficiently small positive $h$, and $\kappaK^{-1}\ll_n1$.
\end{lemma}

\begin{proof}
Put $L=\log D$. Stark's Lemma~3 gives at most one zero in
$\operatorname{Re}s\ge1-(4L)^{-1}$,
$|\operatorname{Im}s|\le(4L)^{-1}$; such a zero is real and simple.
His Lemma~8, which also applies to nonnormal fields, says that a real
zero $\beta\ge1-(4n!L)^{-1}$ forces a quadratic subfield of $K$
\cite[Lemmas~3 and~8]{Stark1974}. Since $n$ is odd, there is no such
subfield. Consequently, for a constant $c_n>0$, every nontrivial zero
$\rho$ satisfies $|1-\rho|\ge c_n/L$.

The real logarithmic derivative of the completed zeta function gives
\[
 S_K(\sigma):=\sum_\rho\operatorname{Re}\frac1{\sigma-\rho}
 =\frac1\sigma+\frac1{\sigma-1}+\frac L2
   +\frac{\zeta_K'}{\zeta_K}(\sigma)+G_n(\sigma),\qquad \sigma>1,
\]
where zeros are counted with multiplicity and
\[
 G_n(\sigma)=\frac n2
 \left(\frac{\Gamma'(\sigma/2)}{\Gamma(\sigma/2)}-\log\pi\right).
\]
This is the Hadamard identity in
\cite[Equation~(9)]{Stark1974}. The zero terms are nonnegative,
$G_n$ is bounded on $[1,2]$, and the Euler product gives
$\zeta_K'/\zeta_K\le0$ there. Thus, at $\sigma_0=1+L^{-1}$,
$S_K(\sigma_0)\ll_n L$.

This bound persists uniformly for $1<\sigma\le1+4/L$.
Indeed, for $\rho=\beta+i\gamma$, set
$X=L(1-\beta)$, $Y=L\gamma$, and $U=L(\sigma-1)$.
Then $X\ge0$, $0<U\le4$, and $|X+iY|\ge c_n$. Hence
\[
 \frac{\operatorname{Re}(\sigma-\rho)^{-1}}
      {\operatorname{Re}(\sigma_0-\rho)^{-1}}
 =\frac{X+U}{X+1}
   \frac{(X+1)^2+Y^2}{(X+U)^2+Y^2}
 \le4(1+3/c_n)^2.
\]
Summing gives $0\le S_K(\sigma)\ll_n L$.
For $F_K(\sigma)=(\sigma-1)\zeta_K(\sigma)$, it follows that
\[
 \frac{F_K'}{F_K}(\sigma)
 =S_K(\sigma)-\frac1\sigma-\frac L2-G_n(\sigma)
 =O_n(L).
\]
Since $F_K$ extends analytically to $1$ with $F_K(1)=\kappaK>0$,
integration over an interval of length at most $4/L$ proves the
comparison in \eqref{eq:residue}. At $\sigma_0$, the inequality
$\zeta_K(\sigma_0)\ge1$ also gives $\kappaK\gg_n L^{-1}$.

For completeness, only finitely many fields of this fixed degree
have discriminant below the cutoff. One can see this directly here:
nonzero vectors of $D^{-1/(2n)}\sigma(R)$ have absolute coordinate
product at least $D^{-1/2}$, so Mahler's criterion puts these lattices
in a fixed compact set. Local basis lifts and a finite cover give
bases of $R$ with uniformly bounded conjugates. A basis element
outside $\Q$ generates $K$, since $n$ is prime; its monic integral
minimal polynomial therefore has bounded coefficients. There are
only finitely many such polynomials. The simple pole of each of
these finitely many zeta functions now gives the remaining claims
uniformly for sufficiently small $h>0$.
\end{proof}

We also use the following uniform form of Shiu's theorem
\cite{Shiu1980,Pollack2020}: if $b$ is nonnegative and multiplicative
and $b(m)\le d_n(m)$, where $d_n$ is the $n$-fold divisor function,
then
\begin{equation}\label{eq:shiu}
 \sum_{u-v<m\le u}b(m)
 \ll_n\frac v{\log u}
       \exp\left(\sum_{\ell\le u}\frac{b(\ell)}\ell\right)
 \qquad(u^{1/2}\le v\le u).
\end{equation}
Here $\ell$ denotes a rational prime, and the threshold for $u$
depends only on $n$. To check uniformity, the growth conditions in
\cite[Theorem~1.1]{Pollack2020}, with all excluded residue sets
empty, are supplied by
$d_n(\ell^j)=\binom{j+n-1}{n-1}\le n^j$ and
$d_n(m)\ll_{n,\eps}m^\eps$. Its strict interval condition
$u^\beta<v\le u$ is satisfied by taking $\beta=1/4$; thus
\eqref{eq:shiu} includes both endpoints of the displayed range.

\begin{proposition}\label{prop:weighted-count}
Fix $A_0\ge1$ and $0<\eta\le1$. For $Q$ sufficiently large in terms
of $n,A_0,\eta$,
\begin{equation}\label{eq:weighted-count}
 \sum_{x-y<\N\mathfrak a\le x}P(\mathfrak a)
 \ll_n \frac{\kappaK y}{q}
 \qquad
 (Q^{3/4}\le x\le A_0Q,\quad \eta x\le y\le x).
\end{equation}
The implied constant is independent of $A_0$ and $\eta$.
\end{proposition}

\begin{proof}
Set $b(m)=a_K(m)\one_{(m,q)=1}$. This function is nonnegative,
multiplicative, and bounded by $d_n$. Indeed, at every prime $\ell$,
including the ramified primes, the local generating series for
$a_K$ is $\prod_{\mathfrak p\mid\ell}(1-T^{f_{\mathfrak p}})^{-1}$;
it is coefficientwise bounded by $(1-T)^{-n}$.
Decomposing an ideal into its part over $S$ and its coprime part,
the sum in \eqref{eq:weighted-count} becomes
\begin{equation}\label{eq:weighted-split}
 \sum_{\mathbf k}P(\mathbf k)
       \sum_{(x-y)/d(\mathbf k)<m\le x/d(\mathbf k)}b(m).
\end{equation}

First suppose $d=d(\mathbf k)\le Q^{1/8}$, and put $u=x/d$.
Then $u\ge Q^{5/8}\ge D^{1/4}$. For $h=1/\log u$,
Lemma~\ref{lem:residue} gives
$\zeta_K(1+h)\ll_n\kappaK\log u$.
There is no upper restriction on $u$: increasing $u$ only decreases
$h$, so the lemma applies even when $D$ stays bounded and $q$ grows.
Positivity of the Euler logarithm yields
\[
 \sum_{\ell\le u}a_K(\ell)\ell^{-1-h}\le\log\zeta_K(1+h).
\]
Since $a_K(\ell)\le n$ and
$\sum_{\ell\le u}(\log\ell)/\ell\ll\log u$, replacing
$\ell^{-1-h}$ by $\ell^{-1}$ costs only $O_n(1)$.
The latter prime-sum bound follows, for example, by comparing
$\sum_{\ell\le u}\lfloor u/\ell\rfloor\log\ell$ with
$\log(\lfloor u\rfloor!)$. Therefore
\[
 \exp\left(\sum_{\ell\le u}\frac{a_K(\ell)}\ell\right)
 \ll_n\kappaK\log u.
\]
Furthermore,
\[
 \log Z_S=\sum_{\ell\in S}\frac{a_K(\ell)}\ell+O_n(1),
 \qquad
 \sum_{\substack{\ell\in S\\\ell>u}}
       \frac{a_K(\ell)}\ell
 \ll_n\frac{\log(2q)}{Q^{5/8}}\ll_n1.
\]
The first error is bounded by a constant times
$\sum_\ell\ell^{-2}$. Removing the primes of $S$ thus gives
\begin{equation}\label{eq:coprime-euler}
 \exp\left(\sum_{\ell\le u}\frac{b(\ell)}\ell\right)
 \ll_n\frac{\kappaK\log u}{Z_S}.
\end{equation}
For $Q$ sufficiently large at fixed $\eta$, the length
$v=y/d\ge\eta u$ is at least $u^{1/2}$.
Applying \eqref{eq:shiu}, then \eqref{eq:coprime-euler} and
\eqref{eq:shell-sum}, bounds the contribution of these $d$ by
\[
 C_n\frac{\kappaK y}{Z_S}
       \sum_{d(\mathbf k)\le Q^{1/8}}\frac{P(\mathbf k)}{d(\mathbf k)}
 \le C_n\frac{\kappaK y}{q}.
\]

For $d>Q^{1/8}$, terms with $d>x$ vanish because the inner sum
in \eqref{eq:weighted-split} then contains no positive integer.
For the remaining terms, $x/d\ge1$, and the elementary bound
$\sum_{m\le z}d_n(m)\le z(1+\log z)^{n-1}$ for $z\ge1$
follows by summing reciprocals in the first $n-1$ factors.
Using \eqref{eq:weight-tail}, with $s=1/(2n)$, the remaining
contribution to \eqref{eq:weighted-split} is at most
\[
 C_n\frac{x}{q}Q^{-s/16}(1+\log x)^{n-1}
 =o\left(\frac{\kappaK x}{q}\right).
\]
This is uniform in the stated range: Lemma~\ref{lem:residue}
gives $\kappaK^{-1}\ll_n\log(3D)$, and $D\le Q^2$.
At fixed $A_0,\eta$, increase the threshold until this term is at
most $\kappaK\eta x/q\le\kappaK y/q$. The implied constant has
not changed with $\eta$, proving the proposition.
\end{proof}
